\subsection{Задача о наибольшей общей подпоследовательности}

\textbf{Определение 1.1} (подпоследовательность).
\emph{Подпоследовательностью} строки $A = a_1 a_2 \ldots a_n$ называется строка вида $a_{i_1} a_{i_2} \ldots a_{i_k}$, где $1 \leq i_1 < i_2 < \ldots < i_k \leq n$.

\textbf{Определение 1.2} (LCS).
Для двух строк $A$ и $B$ их \emph{наибольшая общая подпоследовательность} (LCS) --- это подпоследовательность, являющаяся одновременно подпоследовательностью $A$ и $B$ и имеющая максимально возможную длину. Длина LCS обозначается $\mathrm{lcs}(A, B)$. Заметим, что сама LCS может быть не единственной, однако её длина $\mathrm{lcs}(A, B)$ определена однозначно.

Классический ДП-алгоритм Вагнера--Фишера~\cite{wagner1974} вычисляет $\mathrm{lcs}(A, B)$ за время $O(|A| \cdot |B|)$ и память $O(\min(|A|, |B|))$ по рекуррентному соотношению
\[
    f[i][j] =
    \begin{cases}
        0, & i = 0 \text{ или } j = 0,\\
        f[i-1][j-1] + 1, & a_i = b_j,\\
        \max(f[i-1][j],\, f[i][j-1]), & a_i \neq b_j.
    \end{cases}
\]
Корректность рекуррентного соотношения следует из оптимальной подструктуры: при $a_i = b_j$ оптимальное решение пары $(A[1..i], B[1..j])$ продолжает LCS пары $(A[1..i-1], B[1..j-1])$ символом $a_i = b_j$; при $a_i \neq b_j$ хотя бы один из последних символов не входит в LCS, что даёт максимум по двум подзадачам. Оценка памяти $O(\min(|A|, |B|))$ достигается за счёт хранения только двух соседних строк ДП-таблицы.

\subsection{Случайная модель}

Пусть $\sigma \geq 2$ --- размер алфавита, обозначим алфавит $\Sigma_\sigma = \{0, 1, \ldots, \sigma-1\}$. \emph{Случайной моделью LCS} называется выбор двух независимых строк $A, B \in \Sigma_\sigma^n$, в которых каждый символ независимо и равновероятно выбирается из $\Sigma_\sigma$. Длина LCS становится при этом случайной величиной, обозначаемой $L_n = L_n^{(\sigma)}$.

\subsection{Константа Чватала--Санкоффа}

\textbf{Теорема 1.1} (Чватал, Санкофф, 1975~\cite{chvatal1975}).
\emph{Существует предел}
\[
    \gamma_\sigma := \lim_{n \to \infty} \frac{\mathbb{E}[L_n^{(\sigma)}]}{n} \in [0, 1].
\]

Существование предела следует из \emph{супераддитивности} математического ожидания: $\mathbb{E}[L_{n+m}] \geq \mathbb{E}[L_n] + \mathbb{E}[L_m]$. Действительно, если в первых $n$ символах найдена общая подпоследовательность длины $\ell_1$, а во вторых $m$ независимо --- длины $\ell_2$, их конкатенация даёт общую подпоследовательность пары полных строк длины $\ell_1 + \ell_2$, что и означает супераддитивность по математическому ожиданию. Применение леммы Фекете о супераддитивных последовательностях~\cite{fekete1923} гарантирует существование предела $\lim_{n \to \infty} \mathbb{E}[L_n]/n$. Точное значение $\gamma_\sigma$ не известно ни при каком $\sigma$.

Стиил~\cite{steele1986} установил, что $L_n$ сильно концентрируется около своего среднего: $\mathbb{D}[L_n] = O(n)$, откуда $L_n / n \to \gamma_\sigma$ почти наверное. Иными словами, длина LCS асимптотически линейна с коэффициентом $\gamma_\sigma$.

\subsection{Связь с задачей Улама}

Задача Улама о длине наибольшей возрастающей подпоследовательности случайной перестановки длины $n$ имеет известную асимптотику $2\sqrt{n}$~\cite{vershik1977,loganshepp1977,baikdeiftjohansson1999}. LCS двух случайных строк --- естественное обобщение задачи Улама в дискретный алфавит. Асимптотика $\gamma_\sigma$ при $\sigma \to \infty$ имеет вид~\cite{kiwi2005}
\[
    \gamma_\sigma \sim \frac{2}{\sqrt{\sigma}}, \quad \sigma \to \infty.
\]
Интуитивно, при $\sigma \to \infty$ большинство символов в двух случайных строках различны, и задача поиска LCS приближается к задаче нахождения наибольшей возрастающей подпоследовательности после идентификации совпадающих пар индексов, что и объясняет асимптотику $2/\sqrt{\sigma}$.

\subsection{Хронология нижних оценок $\gamma_2$}

В таблице~\ref{tab:history-gamma2} приведена эволюция известных нижних оценок $\gamma_2$ за полвека.

\begin{table}[h!]
\caption{Хронология строгих нижних оценок $\gamma_2$}
\label{tab:history-gamma2}
\centering
\begin{tabular}{|c|l|c|l|}
\hline
Год  & Авторы                            & $\gamma_2 \geq$    & Метод \\
\hline
1975 & Chv\'atal, Sankoff~\cite{chvatal1975}      & $0{,}7615$  & Жадная стратегия \\
1995 & Dan\v{c}\'\i k, Paterson~\cite{dancik1995} & $0{,}7739$  & Конечные автоматы \\
2003 & Lueker~\cite{lueker2003}                   & $0{,}7881$  & Двойственная LP-схема \\
2009 & Lueker~\cite{lueker2009}                   & $0{,}788$   & Обзор и улучшение \\
2024 & Heineman et al.~\cite{heineman2024}        & $0{,}79266$ & Расширенные пространства состояний \\
2026 & \textbf{Данная работа}             & $\mathbf{0{,}7964}$ & \textbf{Оконное ДП} \\
\hline
\end{tabular}
\end{table}

За последние два десятилетия среднегодовой прирост нижней оценки составляет менее $0{,}00025$, что свидетельствует о сложности задачи и о необходимости принципиально новых вычислительных схем.

\subsection{Верхние оценки}

Лучшая известная верхняя оценка для бинарного алфавита --- $\gamma_2 \leq 0{,}826280$ (Lueker, 2009~\cite{lueker2009}). Метод Lueker'а основан на двойственной формулировке Данчика--Патерсона, усиленной компьютерным перебором. Разрыв между лучшей нижней и верхней оценкой $\gamma_2$ составляет $\approx 0{,}034$.

\subsection{Эмпирические оценки}

Численные эксперименты методом Monte-Carlo (классический ДП-LCS на парах независимых строк длины $n = 10\,000{-}20\,000$, усреднение по $500$ реализациям, стандартная ошибка среднего порядка $10^{-3}$) дают следующие приближённые значения константы:
\begin{itemize}
    \item $\gamma_2 \approx 0{,}811$,
    \item $\gamma_3 \approx 0{,}716$,
    \item $\gamma_4 \approx 0{,}653$,
    \item $\gamma_5 \approx 0{,}606$,
\end{itemize}
с дальнейшим монотонным убыванием в соответствии с асимптотикой $\gamma_\sigma \sim 2/\sqrt{\sigma}$. Эти значения служат верхним ориентиром для качества строгих нижних оценок: ни одна нижняя оценка не может превосходить эмпирического среднего, и чем меньше зазор, тем <<качественнее>> метод. Полный набор оценок и параметров Monte-Carlo приведён в~\S5.3.

\subsection{Программные инструменты и предшествующие реализации}

Авторы работы~\cite{heineman2024} опубликовали свой код и данные на GitHub: \url{https://github.com/Statistics-of-Subsequences/PaperMaterials}. Это позволяет сравнивать численные результаты на единых параметрах.

Исходные коды настоящей работы (реализация оконного ДП, универсальная версия для произвольного $\sigma$, Monte-Carlo верификация) доступны в открытом репозитории: \url{https://github.com/Youka4/HSEDiploma}.

\subsection*{Выводы по разделу}
\addcontentsline{toc}{subsection}{Выводы по разделу 1}

\begin{enumerate}
    \item Задача оценки константы Чватала--Санкоффа $\gamma_\sigma$ --- классическая открытая проблема комбинаторики на словах, известная с 1975~года.
    \item Лучшие строгие оценки для бинарного алфавита: $0{,}7927 \leq \gamma_2 \leq 0{,}8263$; эмпирически $\gamma_2 \approx 0{,}811$. Аналогичные разрывы имеют место для $\sigma \geq 3$.
    \item За 2003--2024 годы суммарный прирост нижней оценки $\gamma_2$ составил всего $\approx 0{,}005$. Существующие методы (конечные автоматы, расширенные пространства состояний) упираются в их экспоненциальный рост.
    \item Для дальнейшего продвижения требуется новый класс вычислительных схем, дающий улучшение порядка $10^{-3}$ или более.
\end{enumerate}
